\clearpage\thispagestyle{empty}
\pdfbookmark[0]{Quelle zu Vesta, Ceres, Komet (10. November 1808)}{vestacerescomet1808nov-sources}
\begin{center}\textbf{Quelle zu Vesta, Ceres, Komet (10. November 1808)}\end{center}
\noindent Diese KI-Übertragung ins heutige Deutsch folgt der Mitteilung vom 10. November 1808 in Gauß’ \href{https://gdz.sub.uni-goettingen.de/id/PPN236007467}{\textit{Werke}, Bd.~VI} (Göttingen, 1874), S.~308--309. Das erste dieser Blätter trägt die gedruckte Zahl 380. Der Erstdruck erschien unter \textit{Göttingen} in den \textit{Göttingischen gelehrten Anzeigen}, Stück 180, S.~\href{https://manifests.sub.uni-goettingen.de/iiif/image/gdz:PPN319721507_1808:00001799/full/1600,/0/default.jpg}{1793}, \href{https://manifests.sub.uni-goettingen.de/iiif/image/gdz:PPN319721507_1808:00001800/full/1600,/0/default.jpg}{1794} und \href{https://manifests.sub.uni-goettingen.de/iiif/image/gdz:PPN319721507_1808:00001801/full/1600,/0/default.jpg}{1795}.
\par\medskip
Die Mitteilung unterscheidet Bodes Berliner Beobachtungen, die von ihm mitgeteilten Petersburger Beobachtungen von Schubert und Wisniewsky und die Göttinger Beobachtungen mit Tiarks’ Vergleich. Das Feld für die Deklination vom 2. Juli ist in beiden Drucken leer. Die Deklination vom 22. Juni bleibt in der Mitteilung und im Brief vom 27. Juni jeweils nach der Werkausgabe erhalten. Auch die Schreibung „Nordliche“ in der Kometentabelle folgt der Werkausgabe. KI-Übertragung und KI-Satz.
\clearpage
